\chapter{Manifolds, tangent bundles and vector fields}
\label{chap:manifolds}

In this chapter we set up the basic objects of the library: manifolds modeled on locally convex
spaces, $\Ck{k}$ maps between them, the tangent bundle and the tangent map, vector fields and
their Lie bracket, the double tangent bundle with its canonical flip, and curves with their
velocities. The calculus is the Michal--Bastiani calculus of \cite{LCA}; the references for the
manifold theory are \cite[\S I.3]{Nee06} and \cite[Ch.~3]{GN26}, and \cite{Lan99} for the
Banach case. Throughout, $E$, $F$, $G$ denote real Hausdorff locally convex spaces. We stress
from the start that \emph{no completeness is assumed}: nothing in this chapter needs it.

\section{Manifolds}

\begin{definition}[The Michal--Bastiani groupoid]
  \label{def:mb-groupoid}
  \lean{LCManifolds.mbGroupoid}
  \leanok
  The groupoid $\mathcal{G}_{\mb}(E)$ consists of the homeomorphisms $\varphi\colon U \to V$
  between open subsets of $E$ such that $\varphi$ is smooth on $U$ and $\varphi^{-1}$ is smooth
  on $V$, in the sense of Michal--Bastiani.
\end{definition}

That this is a structure groupoid (closed under composition, inversion and restriction, and
local in nature) follows from the chain rule and the locality of smoothness.

\begin{definition}[Manifolds]
  \label{def:mb-manifold}
  \lean{LCManifolds.MBManifold, LCManifolds.BanachManifold, LCManifolds.HilbertManifold}
  \leanok
  \uses{def:mb-groupoid}
  A \emph{manifold modeled on $E$} is a topological space $M$ with an atlas of charts with
  values in $E$ whose transition maps belong to $\mathcal{G}_{\mb}(E)$
  \cite[Definition I.3.1]{Nee06}. A \emph{Banach manifold} is a manifold modeled on a normed
  space which is complete; a \emph{Hilbert manifold} is one modeled on a complete inner product
  space.
\end{definition}

In Lean, a manifold is a Mathlib \texttt{ChartedSpace} $E$ $M$ with \texttt{HasGroupoid} for
$\mathcal{G}_{\mb}(E)$; Mathlib's charted spaces accept an arbitrary topological model space, so
no new infrastructure is needed. Every chart comes with a \emph{preferred chart}
$c_x = \texttt{chartAt}\ x$ at each point $x \in M$, and we use it systematically.

\begin{remark}[The model space need not be complete]
  \label{rem:man-complete}
  We follow the convention of Glöckner--Neeb \cite[\S5.7]{GN26}: the model space of a manifold
  is not assumed to be complete. The formal theory of this and the following chapters (charts,
  smooth maps, tangent bundle, vector fields, connections, submanifolds) does not use
  completeness, since the calculus of \cite{LCA} replaces integrals in the model space by
  Hahn--Banach arguments. Results that need completeness (integrals with values in $E$,
  ordinary differential equations, exponential maps, inverse function theorems) say so in their
  hypotheses. In particular, a manifold modeled on an incomplete normed space is a manifold in
  our sense, but it is not a Banach manifold. Without completeness things do go wrong:
  \cite[Example 5.6.8]{GN26} is a Lie group modeled on an incomplete normed algebra that has no
  exponential function.
\end{remark}

On normed spaces, Michal--Bastiani smoothness agrees with Fréchet smoothness, so for complete
normed $E$ a Banach manifold is the classical object of \cite{Lan99}.

\begin{proposition}[Products]
  \label{prop:man-product}
  \lean{LCManifolds.prod_mem_mbGroupoid, LCManifolds.SmoothMBOn.prodMap,
    LCManifolds.smoothMBM_fst, LCManifolds.smoothMBM_snd, LCManifolds.ContMDiffMBOn.prodMk}
  \leanok
  \uses{def:mb-manifold, def:contmdiff}
  If $M$ and $N$ are manifolds modeled on $E$ and $F$, then $M \times N$, with the product
  charts, is a manifold modeled on $E \times F$. The projections onto $M$ and $N$ are smooth, and
  a pair $(f, g)$ of $\Ck{k}$ maps $P \to M$, $P \to N$ is a $\Ck{k}$ map $P \to M \times N$.
\end{proposition}

\begin{proof}
  \leanok
  The transition maps of the product atlas are products $\varphi \times \psi$ of transition
  maps, and a product of smooth diffeomorphisms between open sets is again one, since
  $(\varphi \times \psi) = (\varphi \circ \pi_1, \psi \circ \pi_2)$ and pairs of smooth maps are
  smooth. In charts the projections are the linear maps $\pi_1, \pi_2$.
\end{proof}

In Lean the model space of the product is Mathlib's type synonym \texttt{ModelProd} $E$ $F$ of
$E \times F$, which avoids instance diamonds when $M$ or $N$ is itself a product.

\section{Smooth maps}

For maps $f \colon M \to N$, open $U \subseteq M$ and charts $e$ of $M$, $e'$ of $N$, we denote
by
\[
  D(f, U, e, e') = e(\operatorname{dom} e) \cap e^{-1}\bigl(U \cap f^{-1}(\operatorname{dom}
  e')\bigr) \subseteq E
\]
the domain of the local representative $e' \circ f \circ e^{-1}$. It is open when $f$ is
continuous on $U$ (\texttt{LCManifolds.isOpen\_chartDomain}).

\begin{definition}[$\Ck{k}$ maps between manifolds]
  \label{def:contmdiff}
  \lean{LCManifolds.chartDomain, LCManifolds.ContMDiffMBOn, LCManifolds.ContMDiffMB,
    LCManifolds.SmoothMBMOn, LCManifolds.SmoothMBM}
  \leanok
  \uses{def:mb-manifold}
  A map $f \colon M \to N$ is $\Ck{k}$ on an open set $U \subseteq M$ if it is continuous on $U$
  and, for every $x \in U$, the local representative $c'_{f(x)} \circ f \circ c_x^{-1}$ in the
  preferred charts is $\Ck{k}$ on $D(f, U, c_x, c'_{f(x)})$. It is \emph{smooth} on $U$ if it is
  $\Ck{k}$ on $U$ for every $k$.
\end{definition}

Only the preferred charts appear in the definition. The next result says that this is no
restriction.

\begin{theorem}[Independence of the charts]
  \label{thm:man-contmdiff-chart}
  \lean{LCManifolds.smoothMBOn_transition, LCManifolds.smoothMBOn_transition_of_mem_maximalAtlas,
    LCManifolds.ContMDiffMBOn.contDiffMBOn_chart,
    LCManifolds.ContMDiffMBOn.contDiffMBOn_chart_of_mem_maximalAtlas,
    LCManifolds.contMDiffMB_of_chartAt}
  \leanok
  \uses{def:contmdiff}
  If $f$ is $\Ck{k}$ on $U$, then its local representatives $e' \circ f \circ e^{-1}$ in arbitrary
  charts of the atlases (or of the maximal atlases) are $\Ck{k}$ on $D(f, U, e, e')$.
  Conversely, a continuous map whose local representatives in the preferred charts are
  $\Ck{k}$ near the center of each chart is $\Ck{k}$.
\end{theorem}

\begin{proof}
  \leanok
  The transition maps are smooth by definition of the atlas. Near a point $y = e(p)$ we write
  $e' f e^{-1} = (e' c'^{-1}) \circ (c' f c^{-1}) \circ (c e^{-1})$ with $c = c_p$,
  $c' = c'_{f(p)}$; the three factors are $\Ck{k}$ and the chain rule applies. Since $\Ck{k}$ is a
  local property, this suffices.
\end{proof}

\begin{theorem}[Composition]
  \label{thm:man-comp}
  \lean{LCManifolds.ContMDiffMBOn.comp, LCManifolds.SmoothMBMOn.comp, LCManifolds.SmoothMBM.comp,
    LCManifolds.smoothMBMOn_id, LCManifolds.smoothMBMOn_const}
  \leanok
  \uses{def:contmdiff, thm:man-contmdiff-chart}
  If $f$ is $\Ck{k}$ on $U$, $g$ is $\Ck{k}$ on $V$ and $f(U) \subseteq V$, then $g \circ f$ is
  $\Ck{k}$ on $U$. Identities and constant maps are smooth.
\end{theorem}

\begin{proof}
  \leanok
  In the preferred charts at $x$ and $g(f(x))$, near a point we insert the preferred chart $d$ at
  the intermediate point and write the local representative of $g \circ f$ as the composition of
  the local representatives of $g$ and $f$ in the charts involving $d$. These are $\Ck{k}$ by
  Theorem~\ref{thm:man-contmdiff-chart}, and we conclude by the chain rule.
\end{proof}

\begin{proposition}[Consistency with the calculus]
  \label{prop:man-consistency}
  \lean{LCManifolds.contMDiffMBOn_iff_contDiffMBOn}
  \leanok
  \uses{def:contmdiff}
  For a map $f \colon E \to F$ between the model spaces themselves (with their identity charts)
  and $U \subseteq E$ open, $f$ is $\Ck{k}$ on $U$ as a map between manifolds if and only if it
  is $\Ck{k}$ on $U$ in the sense of Michal--Bastiani.
\end{proposition}

\begin{proof}
  \leanok
  The preferred charts are identities, so $D(f, U, \id, \id) = U$ and the local representative
  is $f$ itself.
\end{proof}

\begin{remark}[Universes]
  \label{rem:man-universes}
  The model spaces $E$, $F$, $G$ live in a common universe, because the Michal--Bastiani
  derivative of a map defined on $E$ is a map defined on $E \times E$. The manifolds themselves
  are unrestricted. A consequence that will reappear below is that real-valued functions on
  $M$ are not treated as maps between manifolds (the model space $\R$ need not be in the
  universe of $E$), but through their local representatives.
\end{remark}

\section{The tangent bundle and tangent maps}

\begin{definition}[The tangent bundle]
  \label{def:tangent-bundle}
  \lean{LCManifolds.tangentCoordChange, LCManifolds.tangentBundleCore,
    LCManifolds.TangentBundle, LCManifolds.tangentChart}
  \leanok
  \uses{def:mb-manifold}
  For charts $i$, $j$ of the atlas of $M$ and $x$ in both domains, the \emph{coordinate change}
  of tangent vectors is $v \mapsto \dd(j \circ i^{-1})(i(x))\,v$. The \emph{tangent bundle}
  $\T M$ is the fiber bundle with fiber $E$ obtained by gluing the trivial bundles over the
  chart domains along these coordinate changes. For a chart $i$, the map
  $\T M \supseteq \T M|_{\operatorname{dom} i} \to E \times E$ that sends $(x, v)$ to
  $(i(x), v_i)$, with $v_i$ the coordinates of $v$ in the chart $i$, is the \emph{chart of $\T M$
  over $i$}.
\end{definition}

In Lean, $\T M$ is the total space of a Mathlib \texttt{FiberBundleCore} indexed by the charts of
the atlas, whose index at $x$ is the preferred chart $c_x$. Thus a point of $\T M$ is a pair
$(x, v)$ with $x \in M$ and $v \in E$, and \emph{$v$ is the tangent vector written in the
coordinates of the preferred chart at $x$}. In any other chart $i$, its coordinates are
$\dd(i \circ c_x^{-1})(c_x(x))\,v$. This is a choice of representative, not a restriction:
every construction below (tangent maps, local representatives of vector fields, Lie brackets,
derivations, velocities) is defined in the preferred chart and then shown to transform correctly
in every chart of the atlas. The coordinate changes satisfy the cocycle condition by the chain
rule (\texttt{LCManifolds.tangentCoordChange\_comp}), and they are continuous in $(x, v)$ since
the transition maps are $\Ck{1}$.

\begin{theorem}[$\T M$ is a manifold]
  \label{thm:man-tangent-manifold}
  \lean{LCManifolds.smoothMBOn_tangentChart_transition, LCManifolds.tangentChart_transition_mem,
    LCManifolds.tangentChart_mem_atlas, LCManifolds.tangentChart_transition_apply}
  \leanok
  \uses{def:tangent-bundle}
  The charts of $\T M$ over the charts of $M$ form an atlas modeled on $E \times E$ whose
  transition maps are
  \[
    (y, w) \longmapsto \bigl(\tau(y), \dd\tau(y)\,w\bigr), \qquad \tau = j \circ i^{-1}.
  \]
  In particular they are smooth, and $\T M$ is a manifold modeled on $E \times E$.
\end{theorem}

\begin{proof}
  \leanok
  The formula is the cocycle property of the coordinate changes. The map $(y, w) \mapsto
  \dd\tau(y)\,w$ is smooth because $\tau$ is, by the recursive definition of $\Ck{k}$, and pairs of
  smooth maps are smooth. The inverse transition has the same form with $\tau^{-1}$.
\end{proof}

\begin{definition}[The tangent map]
  \label{def:tangent-map}
  \lean{LCManifolds.tangentMap}
  \leanok
  \uses{def:tangent-bundle, def:contmdiff}
  Let $f \colon M \to N$. The \emph{tangent map} $\T f \colon \T M \to \T N$ sends $(x, v)$ to
  \[
    \bigl(f(x),\ \dd(c'_{f(x)} \circ f \circ c_x^{-1})(c_x(x))\,v\bigr),
  \]
  where $c_x$ and $c'_{f(x)}$ are the preferred charts at $x$ and $f(x)$.
\end{definition}

The definition makes sense for every $f$, since $\dd$ is a total function in Lean; it is
meaningful when $f$ is $\Ck{1}$.

\begin{theorem}[The tangent map in charts; chain rule]
  \label{thm:man-tangent-map}
  \lean{LCManifolds.dMB_eq_comp_of_eventuallyEq, LCManifolds.dMB_chart_comp,
    LCManifolds.dMB_comp_chart, LCManifolds.tangentMap_chart, LCManifolds.tangentMap_id,
    LCManifolds.tangentMap_comp}
  \leanok
  \uses{def:tangent-map, thm:man-contmdiff-chart}
  Let $f \colon M \to N$ be $\Ck{1}$. The derivatives of the local representatives of $f$
  transform under changes of charts by the derivatives of the transition maps. Consequently, in
  the charts of $\T M$ and $\T N$ over charts $i$ of $M$ and $j$ of $N$, the tangent map reads
  \[
    (y, w) \longmapsto \bigl(h(y), \dd h(y)\,w\bigr), \qquad h = j \circ f \circ i^{-1}.
  \]
  Moreover $\T\id = \id$ and, for $\Ck{1}$ maps $f$ and $g$, $\T(g \circ f) = \T g \circ \T f$.
\end{theorem}

\begin{proof}
  \leanok
  Near a point, the local representative of $f$ in one pair of charts agrees with the
  composition of its local representative in another pair with transition maps; the chain rule for maps that agree with a composition near the point gives
  the transformation laws. The formula in charts follows by composing these laws with the
  coordinate changes of $\T M$ and $\T N$, and functoriality is the chain rule in the preferred
  charts.
\end{proof}

\begin{theorem}[Smoothness of the tangent map]
  \label{thm:man-tangent-map-smooth}
  \lean{LCManifolds.ContMDiffMB.tangentMap}
  \leanok
  \uses{thm:man-tangent-map, thm:man-tangent-manifold}
  If $f \colon M \to N$ is $\Ck{k+1}$, then $\T f \colon \T M \to \T N$ is $\Ck{k}$.
\end{theorem}

\begin{proof}
  \leanok
  In the charts of the tangent bundles $\T f$ is $(y, w) \mapsto (h(y), \dd h(y)\,w)$, and
  $\dd h$ is $\Ck{k}$ on its domain because $h$ is $\Ck{k+1}$.
\end{proof}

\section{Vector fields and derivations}

\begin{definition}[Vector fields]
  \label{def:vector-field}
  \lean{LCManifolds.IsVectorField, LCManifolds.localRep}
  \leanok
  \uses{def:tangent-bundle}
  A \emph{vector field} on $M$ is a section $X \colon M \to \T M$ of the tangent bundle. Its
  \emph{local representative} in a chart $i$ is $\widetilde{X}_i \colon i(\operatorname{dom} i)
  \to E$, $\widetilde{X}_i(y) = $ the coordinates of $X(i^{-1}y)$ in the chart $i$.
\end{definition}

\begin{proposition}[Local representatives]
  \label{prop:man-localrep}
  \lean{LCManifolds.localRep_transition, LCManifolds.contDiffMBOn_localRep,
    LCManifolds.contMDiffMB_of_localRep, LCManifolds.localRep_chartAt}
  \leanok
  \uses{def:vector-field, thm:man-tangent-manifold}
  On chart overlaps, $\widetilde{X}_j \circ \tau = \dd\tau(\cdot)\,\widetilde{X}_i$ for the
  transition map $\tau = j \circ i^{-1}$. A vector field is $\Ck{k}$ as a map $M \to \T M$ if and
  only if its local representatives are $\Ck{k}$ (one direction in every chart of the atlas, the
  other from the preferred charts).
\end{proposition}

\begin{proof}
  \leanok
  The transformation law is the cocycle property of the coordinate changes. The local
  representative $\widetilde{X}_i$ is the second component of the local representative of $X$ in
  the charts $i$ and $\T i$, which gives both implications.
\end{proof}

For maps $X, Y$ defined on an open subset of $E$, the \emph{local bracket} is
\[
  [X, Y](y) = \dd Y(y)\,X(y) - \dd X(y)\,Y(y).
\]

\begin{lemma}[The local bracket]
  \label{lem:man-local-bracket}
  \lean{LCManifolds.localBracket, LCManifolds.dMB_dMB_comp_apply,
    LCManifolds.localBracket_transition, LCManifolds.localBracket_jacobi,
    LCManifolds.contDiffMBOn_localBracket}
  \leanok
  Let $\tau$ be $\Ck{2}$, and let $X_i, Y_i, X_j, Y_j$ be $\Ck{1}$ with
  $X_j \circ \tau = \dd\tau \cdot X_i$ and $Y_j \circ \tau = \dd\tau \cdot Y_i$ near $y$. Then
  $[X_j, Y_j](\tau y) = \dd\tau(y)\,[X_i, Y_i](y)$. The local bracket satisfies the Jacobi identity
  for $\Ck{2}$ fields, and it is $\Ck{k}$ for $\Ck{k+1}$ fields.
\end{lemma}

\begin{proof}
  \leanok
  Differentiating $y \mapsto \dd\tau(y)\,X_i(y)$ in the direction $Y_i(y)$ gives
  $\dd^2\tau(y)(Y_i, X_i) + \dd\tau(y)\,\dd X_i(y)\,Y_i$. In the difference defining the bracket
  the second derivatives $\dd^2\tau(y)(Y_i, X_i)$ and $\dd^2\tau(y)(X_i, Y_i)$ cancel by the
  symmetry of the second differential. The Jacobi identity is a direct computation with the
  derivative of the local bracket, where again the second derivatives cancel by symmetry.
\end{proof}

\begin{definition}[The Lie bracket]
  \label{def:man-lie-bracket}
  \lean{LCManifolds.lieBracket}
  \leanok
  \uses{def:vector-field, lem:man-local-bracket}
  The \emph{Lie bracket} of the vector fields $X$ and $Y$ is the vector field whose value at $x$
  is the local bracket $[\widetilde{X}, \widetilde{Y}](c_x(x))$ of the local representatives in
  the preferred chart at $x$.
\end{definition}

\begin{theorem}[The bracket is well defined]
  \label{thm:man-lie-bracket}
  \lean{LCManifolds.localRep_lieBracket, LCManifolds.localRep_lieBracket_eventuallyEq,
    LCManifolds.lieBracket_swap, LCManifolds.lieBracket_jacobi,
    LCManifolds.contMDiffMB_lieBracket}
  \leanok
  \uses{def:man-lie-bracket, prop:man-localrep}
  Let $X$, $Y$ be $\Ck{1}$ vector fields. In every chart $j$ of the atlas, the local
  representative of $[X, Y]$ is the local bracket $[\widetilde{X}_j, \widetilde{Y}_j]$. The
  bracket is antisymmetric, it satisfies the Jacobi identity for $\Ck{2}$ fields, and if $X$ and
  $Y$ are $\Ck{k+1}$, then $[X, Y]$ is $\Ck{k}$.
\end{theorem}

\begin{proof}
  \leanok
  \uses{lem:man-local-bracket}
  The local representatives in the preferred chart at $x$ and in $j$ are related by
  Proposition~\ref{prop:man-localrep}, so Lemma~\ref{lem:man-local-bracket} applies. The other
  statements are then read in a single chart.
\end{proof}

\begin{remark}[No flows]
  \label{rem:man-no-flows}
  In finite dimensions, chart independence of the bracket is usually derived from its
  description as the Lie derivative along the flow of $X$. In infinite dimensions vector fields
  need not have flows, so this argument is not available. We prove it directly, as above: the
  second derivatives of the transition maps cancel by the symmetry of second derivatives.
\end{remark}

\begin{theorem}[Related fields]
  \label{thm:man-related}
  \lean{LCManifolds.IsRelated, LCManifolds.localRep_related, LCManifolds.lieBracket_related}
  \leanok
  \uses{thm:man-lie-bracket, def:tangent-map}
  Let $f \colon M \to N$ be $\Ck{2}$. We say that $X$ on $M$ and $X'$ on $N$ are $f$-related if
  $\T f \circ X = X' \circ f$. If $X, Y$ are $f$-related to $X', Y'$ (all of them $\Ck{1}$), then
  $[X, Y]$ is $f$-related to $[X', Y']$ \cite[\S V.1]{Lan99}.
\end{theorem}

\begin{proof}
  \leanok
  In charts, $\widetilde{X}'_j \circ h = \dd h \cdot \widetilde{X}_i$ with $h = j \circ f \circ
  i^{-1}$. This is the hypothesis of Lemma~\ref{lem:man-local-bracket} with $h$ in place of
  $\tau$, whose proof only uses that $\tau$ is $\Ck{2}$.
\end{proof}

We now let vector fields act on functions. As explained in Remark~\ref{rem:man-universes}, a
function $\varphi \colon M \to \R$ is said to be $\Ck{k}$ when all its local representatives
$\varphi \circ i^{-1}$ are $\Ck{k}$ on $i(\operatorname{dom} i)$.

\begin{definition}[Derivative along a vector field]
  \label{def:man-vf-deriv}
  \lean{LCManifolds.IsCkFunction, LCManifolds.vfDeriv, LCManifolds.smulField}
  \leanok
  \uses{def:vector-field}
  For a vector field $X$ and $\varphi \colon M \to \R$, we let
  $X(\varphi)(x) = \dd(\varphi \circ c_x^{-1})(c_x(x))\,\widetilde{X}(c_x(x))$, computed in the
  preferred chart at $x$. The product $\varphi Y$ of a function and a vector field is
  $(\varphi Y)(x) = \varphi(x)\,Y(x)$.
\end{definition}

\begin{theorem}[Vector fields as derivations]
  \label{thm:man-derivation}
  \lean{LCManifolds.vfDeriv_chart, LCManifolds.vfDeriv_mul, LCManifolds.HasDirDerivAt.smul_fun,
    LCManifolds.isVectorField_smulField, LCManifolds.localRep_smulField,
    LCManifolds.lieBracket_smulField}
  \leanok
  \uses{def:man-vf-deriv, thm:man-lie-bracket}
  Let $\varphi, \psi$ be $\Ck{1}$ functions. Then $X(\varphi)(x) = \dd(\varphi \circ
  j^{-1})(j(x))\,\widetilde{X}_j(j(x))$ in every chart $j$ around $x$, and the Leibniz rules
  \[
    X(\varphi\psi) = X(\varphi)\,\psi + \varphi\,X(\psi), \qquad
    [X, \varphi Y] = X(\varphi)\,Y + \varphi\,[X, Y]
  \]
  hold, the second one for $\Ck{1}$ vector fields $X$, $Y$.
\end{theorem}

\begin{proof}
  \leanok
  Chart independence follows from $\widetilde{X}_j \circ \tau = \dd\tau \cdot \widetilde{X}_i$
  and the chain rule. The first Leibniz rule is the product rule in a chart. For the second one,
  $\widetilde{(\varphi Y)}_i = (\varphi \circ i^{-1})\,\widetilde{Y}_i$, and we differentiate this
  product in the local bracket.
\end{proof}

\begin{theorem}[Vector fields are determined by local functions]
  \label{thm:man-determined}
  \lean{LCManifolds.vfDeriv_clm_comp_chart, LCManifolds.eq_of_vfDeriv_clm_comp_chart}
  \leanok
  \uses{thm:man-derivation}
  Let $c$ be a chart around $x$. For $\lambda \in E'$ we have $X(\lambda \circ c)(x) =
  \lambda(\widetilde{X}_c(c(x)))$. Consequently, if $X(\lambda \circ c)(x) = Y(\lambda \circ c)(x)$
  for every continuous linear functional $\lambda$, then $X(x) = Y(x)$.
\end{theorem}

\begin{proof}
  \leanok
  The first statement is Theorem~\ref{thm:man-derivation} in the chart $c$, since $\lambda$ is
  linear. The second follows because continuous linear functionals separate the points of the
  Hausdorff locally convex space $E$ (Hahn--Banach).
\end{proof}

\begin{remark}[No bump functions]
  \label{rem:man-no-bumps}
  Smooth bump functions need not exist on manifolds modeled on locally convex spaces. Hence a vector field is not determined in general by its
  action on \emph{globally} defined functions, and the classical identification of vector fields
  with derivations of $\Ck{\infty}(M)$ is not available. Theorem~\ref{thm:man-determined} is the
  substitute: functions defined on a chart domain suffice.
\end{remark}

\section{The double tangent bundle}

The double tangent bundle is $\T\T M = \T(\T M)$, the tangent bundle of the manifold $\T M$
(Theorem~\ref{thm:man-tangent-manifold}). We follow \cite[Cap.~5]{Lar22} and \cite[\S
IV.3]{Lan99}. A chart $i$ of $M$ gives a chart of $\T M$ and then a chart of $\T\T M$, in which
points are written $((x, v), (u, w))$: $(x, v)$ is the base point in $\T M$, and $(u, w)$ the
tangent vector at it.

\begin{theorem}[Transition maps of $\T\T M$]
  \label{thm:man-double-transition}
  \lean{LCManifolds.tangentAtlas, LCManifolds.dMB_tangentChart_transition,
    LCManifolds.DoubleTangentBundle, LCManifolds.ddChart, LCManifolds.ddChart_transition_apply}
  \leanok
  \uses{thm:man-tangent-manifold}
  Let $\tau = j \circ i^{-1}$. In the charts of $\T\T M$ over $i$ and $j$, the transition map is
  \[
    ((x, v), (u, w)) \longmapsto \bigl((\tau x, \dd\tau_x v),\ (\dd\tau_x u,\ \dd^2\tau_x(u, v)
    + \dd\tau_x w)\bigr).
  \]
\end{theorem}

\begin{proof}
  \leanok
  The transition maps of $\T\T M$ are those of the tangent bundle of $\T M$, namely the
  transition map $\sigma(x, v) = (\tau x, \dd\tau_x v)$ of $\T M$ together with its derivative.
  Differentiating $\sigma$ in the direction $(u, w)$ gives $(\dd\tau_x u, \dd^2\tau_x(u, v) +
  \dd\tau_x w)$, since $\dd\tau_x$ is linear in $v$.
\end{proof}

\begin{theorem}[The canonical flip]
  \label{thm:man-flip}
  \lean{LCManifolds.swapUV, LCManifolds.ddChart_transition_swap, LCManifolds.ddChartAt,
    LCManifolds.flip, LCManifolds.flip_proj_proj, LCManifolds.ddChart_flip,
    LCManifolds.flip_flip}
  \leanok
  \uses{thm:man-double-transition}
  The swap $((x, v), (u, w)) \mapsto ((x, u), (v, w))$ commutes with the transition maps of
  $\T\T M$. Hence the map $\kappa \colon \T\T M \to \T\T M$ defined as the swap in the preferred
  chart reads as the swap in every chart of $\T\T M$ over a chart of $M$. It covers the identity
  of $M$, and $\kappa \circ \kappa = \id$.
\end{theorem}

\begin{proof}
  \leanok
  By Theorem~\ref{thm:man-double-transition}, the claim amounts to $\dd^2\tau_x(u, v) =
  \dd^2\tau_x(v, u)$, the symmetry of the second differential. The rest is immediate from the
  coordinate expression.
\end{proof}

The smoothness of $\kappa$ is not stated in the library.

\section{Curves and velocities}

Curves are maps $\alpha \colon \R \to M$. We treat them, like functions, through their local
representatives, without regarding $\R$ as a manifold.

\begin{definition}[$\Ck{k}$ curves and lifts]
  \label{def:ck-curve}
  \lean{LCManifolds.IsCkCurve, LCManifolds.liftCoord, LCManifolds.IsCkLift}
  \leanok
  \uses{def:mb-manifold, def:tangent-bundle}
  A curve $\alpha \colon \R \to M$ is $\Ck{k}$ if it is continuous and $i \circ \alpha$ is
  $\Ck{k}$ on the open set $\alpha^{-1}(\operatorname{dom} i)$ for every chart $i$ of the atlas. A
  \emph{lift} of $\alpha$ is a curve $\beta \colon \R \to \T M$ with $\beta(t) \in \T_{\alpha(t)}
  M$; its fiber coordinate in the chart $i$ is the second component $\mu_i(t)$ of the chart of
  $\T M$ over $i$ at $\beta(t)$. The lift is $\Ck{k}$ if every $\mu_i$ is $\Ck{k}$ on
  $\alpha^{-1}(\operatorname{dom} i)$.
\end{definition}

\begin{definition}[Velocity]
  \label{def:velocity}
  \lean{LCManifolds.velocity}
  \leanok
  \uses{def:ck-curve}
  The \emph{velocity} of a curve $\alpha$ is the lift $\dot\alpha(t) = (\alpha(t),
  (c_{\alpha(t)} \circ \alpha)'(t))$, the derivative of the local representative in the preferred
  chart at $\alpha(t)$.
\end{definition}

\begin{theorem}[The velocity in charts]
  \label{thm:man-velocity}
  \lean{LCManifolds.dMB_chart_curve, LCManifolds.liftCoord_velocity,
    LCManifolds.isCkLift_velocity, LCManifolds.continuous_velocity}
  \leanok
  \uses{def:velocity, thm:man-tangent-manifold}
  Let $\alpha$ be a $\Ck{1}$ curve. For charts $i$, $j$ around $\alpha(t)$ we have
  $(j \circ \alpha)'(t) = \dd(j \circ i^{-1})(i(\alpha(t)))\,(i \circ \alpha)'(t)$; hence the fiber
  coordinate of $\dot\alpha(t)$ in any chart $j$ is $(j \circ \alpha)'(t)$. The velocity is a
  continuous curve in $\T M$, and if $\alpha$ is $\Ck{k+1}$, then $\dot\alpha$ is a $\Ck{k}$ lift.
\end{theorem}

\begin{proof}
  \leanok
  Near $t$, $j \circ \alpha = (j \circ i^{-1}) \circ (i \circ \alpha)$, and we apply the chain
  rule. The second statement is the first one with $i$ the preferred chart. In the chart of
  $\T M$ over $i$, $\dot\alpha$ reads $s \mapsto (i(\alpha(s)), (i \circ \alpha)'(s))$, which is
  continuous (respectively $\Ck{k}$ in the second component); the charts of $\T M$ are
  homeomorphisms onto their images.
\end{proof}

\begin{proposition}[Germs and the chain rule for velocities]
  \label{prop:man-velocity-chain}
  \lean{LCManifolds.velocity_congr, LCManifolds.tangentMap_velocity}
  \leanok
  \uses{def:velocity, def:tangent-map}
  The velocity $\dot\alpha(t)$ only depends on the germ of $\alpha$ at $t$. If $p \colon N \to
  M$ is $\Ck{1}$ and $\Gamma$ is a $\Ck{1}$ curve in $N$, then $\T p\,(\dot\Gamma(t))$ is the velocity of $p \circ \Gamma$
  at $t$.
\end{proposition}

\begin{proof}
  \leanok
  Directional derivatives only depend on germs. For the chain rule, near $t$ we write
  $c' \circ p \circ \Gamma = (c' \circ p \circ c^{-1}) \circ (c \circ \Gamma)$ in the preferred
  charts $c$ at $\Gamma(t)$ and $c'$ at $p(\Gamma(t))$, and apply the chain rule.
\end{proof}

\begin{remark}
  \label{rem:man-curves-ode}
  The continuity of $\dot\alpha$ in $\T M$ is what makes lengths of curves meaningful for Finsler
  metrics (see Definition~\ref{def:finsler-length-structure}); the chain rule for velocities is used for quotient
  metrics, where curves in $N$ are pushed down along a projection $p$. On the other hand, we do
  not claim the existence of integral curves of vector fields, or of curves with prescribed
  velocity: this would require existence theorems for ordinary differential equations, which
  fail in general locally convex spaces, and in Banach spaces need completeness.
\end{remark}
